% =====================================================================
% Auxiliary-01 — Brazovskii Universality Programme (conditional)
% Status: [DRAFT] [NEEDS_UPDATE] — Wave 6 auxiliary paper
% AUDIT-FLAG: AUDIT-2026-05-02-Wave7-Aux-Epoch-Overclaim
%   Hostile-referee audit (2026-05-02) flagged "belongs rigorously to
%   the Brazovskii universality class" and "mean-field critical
%   exponents apply directly" as over-claims. Title and abstract have
%   been re-cast as a "conditional Brazovskii-universality programme"
%   per the audit recommendation; the §2.3 Theorem 1 statement has
%   been re-tagged as a conditional programme (T6/CONDITIONAL) rather
%   than a closed theorem.
% Compile: pdflatex Auxiliary-01.tex (×2 for refs)
% Canonical archive: Math97, Math97-AddA, Math97-AddB, Math97-AddC
% =====================================================================
% --- TECT canonical preamble (see ../../papers/_shared/tect-preamble.tex)
% =====================================================================
% TECT canonical LaTeX preamble (revtex4-2 / single-column / theorem-safe)
% =====================================================================
% Maintainer: Jusang Lee (jtkor@outlook.com)
% Binding from: 2026-05-06
% Purpose: a single source-of-truth preamble for every TECT paper
% (Paper-00 .. Paper-10 + Auxiliary notes). Designed to (a) eliminate
% font-corruption on pdflatex (T1 fontenc + lmodern), (b) make
% theorem-heavy mathematical-physics content render cleanly in a
% single-column layout (PRD class option, NOT PRL two-column), (c)
% provide consistent theorem numbering across all TECT papers via
% amsthm with a shared counter set, (d) make hyperref + cleveref
% cross-references resolve cleanly with the standard two-pass
% pdflatex compile.
%
% Usage in each Paper-NN.tex:
%   % =====================================================================
% TECT canonical LaTeX preamble (revtex4-2 / single-column / theorem-safe)
% =====================================================================
% Maintainer: Jusang Lee (jtkor@outlook.com)
% Binding from: 2026-05-06
% Purpose: a single source-of-truth preamble for every TECT paper
% (Paper-00 .. Paper-10 + Auxiliary notes). Designed to (a) eliminate
% font-corruption on pdflatex (T1 fontenc + lmodern), (b) make
% theorem-heavy mathematical-physics content render cleanly in a
% single-column layout (PRD class option, NOT PRL two-column), (c)
% provide consistent theorem numbering across all TECT papers via
% amsthm with a shared counter set, (d) make hyperref + cleveref
% cross-references resolve cleanly with the standard two-pass
% pdflatex compile.
%
% Usage in each Paper-NN.tex:
%   % =====================================================================
% TECT canonical LaTeX preamble (revtex4-2 / single-column / theorem-safe)
% =====================================================================
% Maintainer: Jusang Lee (jtkor@outlook.com)
% Binding from: 2026-05-06
% Purpose: a single source-of-truth preamble for every TECT paper
% (Paper-00 .. Paper-10 + Auxiliary notes). Designed to (a) eliminate
% font-corruption on pdflatex (T1 fontenc + lmodern), (b) make
% theorem-heavy mathematical-physics content render cleanly in a
% single-column layout (PRD class option, NOT PRL two-column), (c)
% provide consistent theorem numbering across all TECT papers via
% amsthm with a shared counter set, (d) make hyperref + cleveref
% cross-references resolve cleanly with the standard two-pass
% pdflatex compile.
%
% Usage in each Paper-NN.tex:
%   \input{../_shared/tect-preamble.tex}
%   \begin{document}
%   ...
%   \end{document}
%
% Build (every paper): `pdflatex Paper-NN && pdflatex Paper-NN`
% (the second pass resolves \ref{} / \cref{} forward references).
% A bash/PowerShell wrapper that automates this is at
% Docs/papers/papers/_shared/build-paper.sh / build-paper.ps1.
% =====================================================================

% ---- Document class --------------------------------------------------
% PRD class option of revtex4-2 with [reprint,onecolumn] gives the same
% APS heading / by-line / abstract style as PRL but in a wider single
% column that handles long display equations and theorem environments
% without column-break artefacts. nofootinbib keeps citations out of
% footnotes. The 11pt option restores standard font metrics; with T1
% fontenc + lmodern this gives non-bitmap glyphs.
\documentclass[%
  aps,prd,reprint,onecolumn,nofootinbib,11pt,%
  amsmath,amssymb,floatfix,longbibliography%
]{revtex4-2}

% ---- Encoding & fonts ------------------------------------------------
\usepackage[T1]{fontenc}        % proper 8-bit font encoding (no font corruption)
\usepackage[utf8]{inputenc}     % allow UTF-8 source
\usepackage{lmodern}            % scalable Latin Modern (replaces bitmap CM)
\usepackage{textcomp}           % extra text symbols
\usepackage{microtype}          % subtle kerning improvements

% ---- UTF-8 → LaTeX character mappings --------------------------------
% Maps the Unicode characters that occur in TECT body text (legacy from
% the chat-content auto-archival rule, CLAUDE.md §4) to LaTeX-safe
% commands. Without these declarations, T1+inputenc rejects the chars
% with "Unicode character X not set up for use with LaTeX".
% Adding declarations centrally (here) is preferable to per-file edits.
\DeclareUnicodeCharacter{00A7}{\S}              % §  (section sign)
\DeclareUnicodeCharacter{00B2}{\ensuremath{^{2}}}        % ²
\DeclareUnicodeCharacter{00B3}{\ensuremath{^{3}}}        % ³
\DeclareUnicodeCharacter{00B5}{\ensuremath{\mu}}         % µ (micro)
\DeclareUnicodeCharacter{00B7}{\ensuremath{\cdot}}       % · (middle dot)
\DeclareUnicodeCharacter{00D7}{\ensuremath{\times}}      % ×
\DeclareUnicodeCharacter{00E4}{\"a}             % ä
\DeclareUnicodeCharacter{00E9}{\'e}             % é
\DeclareUnicodeCharacter{00F6}{\"o}             % ö
\DeclareUnicodeCharacter{00FC}{\"u}             % ü
\DeclareUnicodeCharacter{010C}{\v{C}}           % Č
\DeclareUnicodeCharacter{010D}{\v{c}}           % č
\DeclareUnicodeCharacter{0394}{\ensuremath{\Delta}}      % Δ
\DeclareUnicodeCharacter{03A3}{\ensuremath{\Sigma}}      % Σ
\DeclareUnicodeCharacter{03A9}{\ensuremath{\Omega}}      % Ω
\DeclareUnicodeCharacter{03B1}{\ensuremath{\alpha}}      % α
\DeclareUnicodeCharacter{03B2}{\ensuremath{\beta}}       % β
\DeclareUnicodeCharacter{03B3}{\ensuremath{\gamma}}      % γ
\DeclareUnicodeCharacter{03B4}{\ensuremath{\delta}}      % δ
\DeclareUnicodeCharacter{03B5}{\ensuremath{\varepsilon}} % ε
\DeclareUnicodeCharacter{03BA}{\ensuremath{\kappa}}      % κ
\DeclareUnicodeCharacter{03BB}{\ensuremath{\lambda}}     % λ
\DeclareUnicodeCharacter{03BC}{\ensuremath{\mu}}         % μ
\DeclareUnicodeCharacter{03BD}{\ensuremath{\nu}}         % ν
\DeclareUnicodeCharacter{03C0}{\ensuremath{\pi}}         % π
\DeclareUnicodeCharacter{03C1}{\ensuremath{\rho}}        % ρ
\DeclareUnicodeCharacter{03C3}{\ensuremath{\sigma}}      % σ
\DeclareUnicodeCharacter{03C4}{\ensuremath{\tau}}        % τ
\DeclareUnicodeCharacter{03C6}{\ensuremath{\varphi}}     % φ
\DeclareUnicodeCharacter{03C7}{\ensuremath{\chi}}        % χ
\DeclareUnicodeCharacter{03C8}{\ensuremath{\psi}}        % ψ
\DeclareUnicodeCharacter{03C9}{\ensuremath{\omega}}      % ω
\DeclareUnicodeCharacter{2013}{--}              % – (en dash)
\DeclareUnicodeCharacter{2014}{---}             % — (em dash)
\DeclareUnicodeCharacter{2018}{`}               % ‘
\DeclareUnicodeCharacter{2019}{'}               % ’
\DeclareUnicodeCharacter{201C}{``}              % “
\DeclareUnicodeCharacter{201D}{''}              % ”
\DeclareUnicodeCharacter{2070}{\ensuremath{^{0}}}        % ⁰
\DeclareUnicodeCharacter{2071}{\ensuremath{^{i}}}        % ⁱ
\DeclareUnicodeCharacter{2122}{\textsuperscript{TM}}     % ™
\DeclareUnicodeCharacter{2126}{\ensuremath{\Omega}}      % Ω (ohm sign)
\DeclareUnicodeCharacter{210F}{\ensuremath{\hbar}}       % ℏ
\DeclareUnicodeCharacter{2115}{\ensuremath{\mathbb{N}}}  % ℕ
\DeclareUnicodeCharacter{2102}{\ensuremath{\mathbb{C}}}  % ℂ
\DeclareUnicodeCharacter{211D}{\ensuremath{\mathbb{R}}}  % ℝ
\DeclareUnicodeCharacter{2124}{\ensuremath{\mathbb{Z}}}  % ℤ
\DeclareUnicodeCharacter{211A}{\ensuremath{\mathbb{Q}}}  % ℚ
\DeclareUnicodeCharacter{2192}{\ensuremath{\to}}         % →
\DeclareUnicodeCharacter{21D2}{\ensuremath{\Rightarrow}} % ⇒
\DeclareUnicodeCharacter{2200}{\ensuremath{\forall}}     % ∀
\DeclareUnicodeCharacter{2203}{\ensuremath{\exists}}     % ∃
\DeclareUnicodeCharacter{2208}{\ensuremath{\in}}         % ∈
\DeclareUnicodeCharacter{2209}{\ensuremath{\notin}}      % ∉
\DeclareUnicodeCharacter{2227}{\ensuremath{\wedge}}      % ∧
\DeclareUnicodeCharacter{2228}{\ensuremath{\vee}}        % ∨
\DeclareUnicodeCharacter{2229}{\ensuremath{\cap}}        % ∩
\DeclareUnicodeCharacter{222A}{\ensuremath{\cup}}        % ∪
\DeclareUnicodeCharacter{2245}{\ensuremath{\cong}}       % ≅
\DeclareUnicodeCharacter{2248}{\ensuremath{\approx}}     % ≈
\DeclareUnicodeCharacter{2260}{\ensuremath{\neq}}        % ≠
\DeclareUnicodeCharacter{2264}{\ensuremath{\leq}}        % ≤
\DeclareUnicodeCharacter{2265}{\ensuremath{\geq}}        % ≥
\DeclareUnicodeCharacter{22C5}{\ensuremath{\cdot}}       % ⋅
\DeclareUnicodeCharacter{2295}{\ensuremath{\oplus}}      % ⊕
\DeclareUnicodeCharacter{2297}{\ensuremath{\otimes}}     % ⊗

% ---- Mathematics & symbols ------------------------------------------
\usepackage{amsmath,amssymb,bm,mathrsfs,mathtools}
\usepackage{amsthm}             % theorem environments (load BEFORE hyperref)

% ---- Theorem environments (shared counter set, TECT canonical) -------
% All theorems / lemmas / propositions / corollaries / remarks share a
% single counter so cross-references read consistently across a paper.
% Definitions get a separate counter (definitions are numbered in their
% own sequence by editorial convention).
\newtheorem{theorem}{Theorem}
\newtheorem{lemma}[theorem]{Lemma}
\newtheorem{proposition}[theorem]{Proposition}
\newtheorem{corollary}[theorem]{Corollary}

\theoremstyle{remark}
\newtheorem{remark}[theorem]{Remark}
\newtheorem{example}[theorem]{Example}

\theoremstyle{definition}
\newtheorem{definition}[theorem]{Definition}
\newtheorem{conjecture}[theorem]{Conjecture}
\newtheorem{hypothesis}[theorem]{Hypothesis}

% ---- Tables / floats / graphics --------------------------------------
\usepackage{booktabs}           % \toprule / \midrule / \bottomrule
\usepackage{array}
\usepackage{graphicx}
\usepackage{enumitem}           % \begin{itemize}[leftmargin=...] options
\usepackage{hyphenat}           % \hyp{} for breakable hyphens in \texttt
\usepackage{seqsplit}           % \seqsplit{} for long unbreakable strings
\sloppy                         % allow stretchier inter-word spacing to
                                % reduce overfull \hbox warnings on long URLs
                                % and long \texttt tokens (paragraph-level).
\emergencystretch=3em           % last-resort hbox stretch budget

% ---- Cross-references (hyperref last among the standard packages,
%      cleveref AFTER hyperref) ----------------------------------------
\usepackage[%
  colorlinks=true,%
  linkcolor=blue,%
  citecolor=blue,%
  urlcolor=blue,%
  breaklinks=true,%
  bookmarks=true,%
  bookmarksnumbered=true%
]{hyperref}
\usepackage[capitalise,nameinlink]{cleveref}

% ---- TECT-canonical author block macros -----------------------------
% (defined here so every paper has the same author / affiliation style)
\newcommand{\tectauthor}{%
  \author{Jusang Lee}%
  \email{jtkor@outlook.com}%
  \affiliation{Independent researcher, Republic of Korea}%
}

% ---- TECT-canonical archive footer ----------------------------------
\newcommand{\tectarchivefooter}{%
  \par\medskip
  \noindent\small\textit{Canonical archive}:
  \url{https://github.com/TECT-OpenPhysics/TECT}.\par%
}

% ---- TECT TODO / AUDIT-FLAG / HONEST-SCOPE inline marks --------------
% Used in drafts to highlight pending audit items without disturbing
% the body text style. Suppress via \tectsuppressmarks (define empty).
\providecommand{\tectauditflag}[1]{\textcolor{red}{\textbf{[AUDIT-FLAG: #1]}}}
\providecommand{\tecttodo}[1]{\textcolor{orange}{\textbf{[TODO: #1]}}}
\providecommand{\tecthonestscope}[1]{\textit{Honest scope: #1.}}

% =====================================================================
% End of TECT canonical preamble
% =====================================================================

%   \begin{document}
%   ...
%   \end{document}
%
% Build (every paper): `pdflatex Paper-NN && pdflatex Paper-NN`
% (the second pass resolves \ref{} / \cref{} forward references).
% A bash/PowerShell wrapper that automates this is at
% Docs/papers/papers/_shared/build-paper.sh / build-paper.ps1.
% =====================================================================

% ---- Document class --------------------------------------------------
% PRD class option of revtex4-2 with [reprint,onecolumn] gives the same
% APS heading / by-line / abstract style as PRL but in a wider single
% column that handles long display equations and theorem environments
% without column-break artefacts. nofootinbib keeps citations out of
% footnotes. The 11pt option restores standard font metrics; with T1
% fontenc + lmodern this gives non-bitmap glyphs.
\documentclass[%
  aps,prd,reprint,onecolumn,nofootinbib,11pt,%
  amsmath,amssymb,floatfix,longbibliography%
]{revtex4-2}

% ---- Encoding & fonts ------------------------------------------------
\usepackage[T1]{fontenc}        % proper 8-bit font encoding (no font corruption)
\usepackage[utf8]{inputenc}     % allow UTF-8 source
\usepackage{lmodern}            % scalable Latin Modern (replaces bitmap CM)
\usepackage{textcomp}           % extra text symbols
\usepackage{microtype}          % subtle kerning improvements

% ---- UTF-8 → LaTeX character mappings --------------------------------
% Maps the Unicode characters that occur in TECT body text (legacy from
% the chat-content auto-archival rule, CLAUDE.md §4) to LaTeX-safe
% commands. Without these declarations, T1+inputenc rejects the chars
% with "Unicode character X not set up for use with LaTeX".
% Adding declarations centrally (here) is preferable to per-file edits.
\DeclareUnicodeCharacter{00A7}{\S}              % §  (section sign)
\DeclareUnicodeCharacter{00B2}{\ensuremath{^{2}}}        % ²
\DeclareUnicodeCharacter{00B3}{\ensuremath{^{3}}}        % ³
\DeclareUnicodeCharacter{00B5}{\ensuremath{\mu}}         % µ (micro)
\DeclareUnicodeCharacter{00B7}{\ensuremath{\cdot}}       % · (middle dot)
\DeclareUnicodeCharacter{00D7}{\ensuremath{\times}}      % ×
\DeclareUnicodeCharacter{00E4}{\"a}             % ä
\DeclareUnicodeCharacter{00E9}{\'e}             % é
\DeclareUnicodeCharacter{00F6}{\"o}             % ö
\DeclareUnicodeCharacter{00FC}{\"u}             % ü
\DeclareUnicodeCharacter{010C}{\v{C}}           % Č
\DeclareUnicodeCharacter{010D}{\v{c}}           % č
\DeclareUnicodeCharacter{0394}{\ensuremath{\Delta}}      % Δ
\DeclareUnicodeCharacter{03A3}{\ensuremath{\Sigma}}      % Σ
\DeclareUnicodeCharacter{03A9}{\ensuremath{\Omega}}      % Ω
\DeclareUnicodeCharacter{03B1}{\ensuremath{\alpha}}      % α
\DeclareUnicodeCharacter{03B2}{\ensuremath{\beta}}       % β
\DeclareUnicodeCharacter{03B3}{\ensuremath{\gamma}}      % γ
\DeclareUnicodeCharacter{03B4}{\ensuremath{\delta}}      % δ
\DeclareUnicodeCharacter{03B5}{\ensuremath{\varepsilon}} % ε
\DeclareUnicodeCharacter{03BA}{\ensuremath{\kappa}}      % κ
\DeclareUnicodeCharacter{03BB}{\ensuremath{\lambda}}     % λ
\DeclareUnicodeCharacter{03BC}{\ensuremath{\mu}}         % μ
\DeclareUnicodeCharacter{03BD}{\ensuremath{\nu}}         % ν
\DeclareUnicodeCharacter{03C0}{\ensuremath{\pi}}         % π
\DeclareUnicodeCharacter{03C1}{\ensuremath{\rho}}        % ρ
\DeclareUnicodeCharacter{03C3}{\ensuremath{\sigma}}      % σ
\DeclareUnicodeCharacter{03C4}{\ensuremath{\tau}}        % τ
\DeclareUnicodeCharacter{03C6}{\ensuremath{\varphi}}     % φ
\DeclareUnicodeCharacter{03C7}{\ensuremath{\chi}}        % χ
\DeclareUnicodeCharacter{03C8}{\ensuremath{\psi}}        % ψ
\DeclareUnicodeCharacter{03C9}{\ensuremath{\omega}}      % ω
\DeclareUnicodeCharacter{2013}{--}              % – (en dash)
\DeclareUnicodeCharacter{2014}{---}             % — (em dash)
\DeclareUnicodeCharacter{2018}{`}               % ‘
\DeclareUnicodeCharacter{2019}{'}               % ’
\DeclareUnicodeCharacter{201C}{``}              % “
\DeclareUnicodeCharacter{201D}{''}              % ”
\DeclareUnicodeCharacter{2070}{\ensuremath{^{0}}}        % ⁰
\DeclareUnicodeCharacter{2071}{\ensuremath{^{i}}}        % ⁱ
\DeclareUnicodeCharacter{2122}{\textsuperscript{TM}}     % ™
\DeclareUnicodeCharacter{2126}{\ensuremath{\Omega}}      % Ω (ohm sign)
\DeclareUnicodeCharacter{210F}{\ensuremath{\hbar}}       % ℏ
\DeclareUnicodeCharacter{2115}{\ensuremath{\mathbb{N}}}  % ℕ
\DeclareUnicodeCharacter{2102}{\ensuremath{\mathbb{C}}}  % ℂ
\DeclareUnicodeCharacter{211D}{\ensuremath{\mathbb{R}}}  % ℝ
\DeclareUnicodeCharacter{2124}{\ensuremath{\mathbb{Z}}}  % ℤ
\DeclareUnicodeCharacter{211A}{\ensuremath{\mathbb{Q}}}  % ℚ
\DeclareUnicodeCharacter{2192}{\ensuremath{\to}}         % →
\DeclareUnicodeCharacter{21D2}{\ensuremath{\Rightarrow}} % ⇒
\DeclareUnicodeCharacter{2200}{\ensuremath{\forall}}     % ∀
\DeclareUnicodeCharacter{2203}{\ensuremath{\exists}}     % ∃
\DeclareUnicodeCharacter{2208}{\ensuremath{\in}}         % ∈
\DeclareUnicodeCharacter{2209}{\ensuremath{\notin}}      % ∉
\DeclareUnicodeCharacter{2227}{\ensuremath{\wedge}}      % ∧
\DeclareUnicodeCharacter{2228}{\ensuremath{\vee}}        % ∨
\DeclareUnicodeCharacter{2229}{\ensuremath{\cap}}        % ∩
\DeclareUnicodeCharacter{222A}{\ensuremath{\cup}}        % ∪
\DeclareUnicodeCharacter{2245}{\ensuremath{\cong}}       % ≅
\DeclareUnicodeCharacter{2248}{\ensuremath{\approx}}     % ≈
\DeclareUnicodeCharacter{2260}{\ensuremath{\neq}}        % ≠
\DeclareUnicodeCharacter{2264}{\ensuremath{\leq}}        % ≤
\DeclareUnicodeCharacter{2265}{\ensuremath{\geq}}        % ≥
\DeclareUnicodeCharacter{22C5}{\ensuremath{\cdot}}       % ⋅
\DeclareUnicodeCharacter{2295}{\ensuremath{\oplus}}      % ⊕
\DeclareUnicodeCharacter{2297}{\ensuremath{\otimes}}     % ⊗

% ---- Mathematics & symbols ------------------------------------------
\usepackage{amsmath,amssymb,bm,mathrsfs,mathtools}
\usepackage{amsthm}             % theorem environments (load BEFORE hyperref)

% ---- Theorem environments (shared counter set, TECT canonical) -------
% All theorems / lemmas / propositions / corollaries / remarks share a
% single counter so cross-references read consistently across a paper.
% Definitions get a separate counter (definitions are numbered in their
% own sequence by editorial convention).
\newtheorem{theorem}{Theorem}
\newtheorem{lemma}[theorem]{Lemma}
\newtheorem{proposition}[theorem]{Proposition}
\newtheorem{corollary}[theorem]{Corollary}

\theoremstyle{remark}
\newtheorem{remark}[theorem]{Remark}
\newtheorem{example}[theorem]{Example}

\theoremstyle{definition}
\newtheorem{definition}[theorem]{Definition}
\newtheorem{conjecture}[theorem]{Conjecture}
\newtheorem{hypothesis}[theorem]{Hypothesis}

% ---- Tables / floats / graphics --------------------------------------
\usepackage{booktabs}           % \toprule / \midrule / \bottomrule
\usepackage{array}
\usepackage{graphicx}
\usepackage{enumitem}           % \begin{itemize}[leftmargin=...] options
\usepackage{hyphenat}           % \hyp{} for breakable hyphens in \texttt
\usepackage{seqsplit}           % \seqsplit{} for long unbreakable strings
\sloppy                         % allow stretchier inter-word spacing to
                                % reduce overfull \hbox warnings on long URLs
                                % and long \texttt tokens (paragraph-level).
\emergencystretch=3em           % last-resort hbox stretch budget

% ---- Cross-references (hyperref last among the standard packages,
%      cleveref AFTER hyperref) ----------------------------------------
\usepackage[%
  colorlinks=true,%
  linkcolor=blue,%
  citecolor=blue,%
  urlcolor=blue,%
  breaklinks=true,%
  bookmarks=true,%
  bookmarksnumbered=true%
]{hyperref}
\usepackage[capitalise,nameinlink]{cleveref}

% ---- TECT-canonical author block macros -----------------------------
% (defined here so every paper has the same author / affiliation style)
\newcommand{\tectauthor}{%
  \author{Jusang Lee}%
  \email{jtkor@outlook.com}%
  \affiliation{Independent researcher, Republic of Korea}%
}

% ---- TECT-canonical archive footer ----------------------------------
\newcommand{\tectarchivefooter}{%
  \par\medskip
  \noindent\small\textit{Canonical archive}:
  \url{https://github.com/TECT-OpenPhysics/TECT}.\par%
}

% ---- TECT TODO / AUDIT-FLAG / HONEST-SCOPE inline marks --------------
% Used in drafts to highlight pending audit items without disturbing
% the body text style. Suppress via \tectsuppressmarks (define empty).
\providecommand{\tectauditflag}[1]{\textcolor{red}{\textbf{[AUDIT-FLAG: #1]}}}
\providecommand{\tecttodo}[1]{\textcolor{orange}{\textbf{[TODO: #1]}}}
\providecommand{\tecthonestscope}[1]{\textit{Honest scope: #1.}}

% =====================================================================
% End of TECT canonical preamble
% =====================================================================

%   \begin{document}
%   ...
%   \end{document}
%
% Build (every paper): `pdflatex Paper-NN && pdflatex Paper-NN`
% (the second pass resolves \ref{} / \cref{} forward references).
% A bash/PowerShell wrapper that automates this is at
% Docs/papers/papers/_shared/build-paper.sh / build-paper.ps1.
% =====================================================================

% ---- Document class --------------------------------------------------
% PRD class option of revtex4-2 with [reprint,onecolumn] gives the same
% APS heading / by-line / abstract style as PRL but in a wider single
% column that handles long display equations and theorem environments
% without column-break artefacts. nofootinbib keeps citations out of
% footnotes. The 11pt option restores standard font metrics; with T1
% fontenc + lmodern this gives non-bitmap glyphs.
\documentclass[%
  aps,prd,reprint,onecolumn,nofootinbib,11pt,%
  amsmath,amssymb,floatfix,longbibliography%
]{revtex4-2}

% ---- Encoding & fonts ------------------------------------------------
\usepackage[T1]{fontenc}        % proper 8-bit font encoding (no font corruption)
\usepackage[utf8]{inputenc}     % allow UTF-8 source
\usepackage{lmodern}            % scalable Latin Modern (replaces bitmap CM)
\usepackage{textcomp}           % extra text symbols
\usepackage{microtype}          % subtle kerning improvements

% ---- UTF-8 → LaTeX character mappings --------------------------------
% Maps the Unicode characters that occur in TECT body text (legacy from
% the chat-content auto-archival rule, CLAUDE.md §4) to LaTeX-safe
% commands. Without these declarations, T1+inputenc rejects the chars
% with "Unicode character X not set up for use with LaTeX".
% Adding declarations centrally (here) is preferable to per-file edits.
\DeclareUnicodeCharacter{00A7}{\S}              % §  (section sign)
\DeclareUnicodeCharacter{00B2}{\ensuremath{^{2}}}        % ²
\DeclareUnicodeCharacter{00B3}{\ensuremath{^{3}}}        % ³
\DeclareUnicodeCharacter{00B5}{\ensuremath{\mu}}         % µ (micro)
\DeclareUnicodeCharacter{00B7}{\ensuremath{\cdot}}       % · (middle dot)
\DeclareUnicodeCharacter{00D7}{\ensuremath{\times}}      % ×
\DeclareUnicodeCharacter{00E4}{\"a}             % ä
\DeclareUnicodeCharacter{00E9}{\'e}             % é
\DeclareUnicodeCharacter{00F6}{\"o}             % ö
\DeclareUnicodeCharacter{00FC}{\"u}             % ü
\DeclareUnicodeCharacter{010C}{\v{C}}           % Č
\DeclareUnicodeCharacter{010D}{\v{c}}           % č
\DeclareUnicodeCharacter{0394}{\ensuremath{\Delta}}      % Δ
\DeclareUnicodeCharacter{03A3}{\ensuremath{\Sigma}}      % Σ
\DeclareUnicodeCharacter{03A9}{\ensuremath{\Omega}}      % Ω
\DeclareUnicodeCharacter{03B1}{\ensuremath{\alpha}}      % α
\DeclareUnicodeCharacter{03B2}{\ensuremath{\beta}}       % β
\DeclareUnicodeCharacter{03B3}{\ensuremath{\gamma}}      % γ
\DeclareUnicodeCharacter{03B4}{\ensuremath{\delta}}      % δ
\DeclareUnicodeCharacter{03B5}{\ensuremath{\varepsilon}} % ε
\DeclareUnicodeCharacter{03BA}{\ensuremath{\kappa}}      % κ
\DeclareUnicodeCharacter{03BB}{\ensuremath{\lambda}}     % λ
\DeclareUnicodeCharacter{03BC}{\ensuremath{\mu}}         % μ
\DeclareUnicodeCharacter{03BD}{\ensuremath{\nu}}         % ν
\DeclareUnicodeCharacter{03C0}{\ensuremath{\pi}}         % π
\DeclareUnicodeCharacter{03C1}{\ensuremath{\rho}}        % ρ
\DeclareUnicodeCharacter{03C3}{\ensuremath{\sigma}}      % σ
\DeclareUnicodeCharacter{03C4}{\ensuremath{\tau}}        % τ
\DeclareUnicodeCharacter{03C6}{\ensuremath{\varphi}}     % φ
\DeclareUnicodeCharacter{03C7}{\ensuremath{\chi}}        % χ
\DeclareUnicodeCharacter{03C8}{\ensuremath{\psi}}        % ψ
\DeclareUnicodeCharacter{03C9}{\ensuremath{\omega}}      % ω
\DeclareUnicodeCharacter{2013}{--}              % – (en dash)
\DeclareUnicodeCharacter{2014}{---}             % — (em dash)
\DeclareUnicodeCharacter{2018}{`}               % ‘
\DeclareUnicodeCharacter{2019}{'}               % ’
\DeclareUnicodeCharacter{201C}{``}              % “
\DeclareUnicodeCharacter{201D}{''}              % ”
\DeclareUnicodeCharacter{2070}{\ensuremath{^{0}}}        % ⁰
\DeclareUnicodeCharacter{2071}{\ensuremath{^{i}}}        % ⁱ
\DeclareUnicodeCharacter{2122}{\textsuperscript{TM}}     % ™
\DeclareUnicodeCharacter{2126}{\ensuremath{\Omega}}      % Ω (ohm sign)
\DeclareUnicodeCharacter{210F}{\ensuremath{\hbar}}       % ℏ
\DeclareUnicodeCharacter{2115}{\ensuremath{\mathbb{N}}}  % ℕ
\DeclareUnicodeCharacter{2102}{\ensuremath{\mathbb{C}}}  % ℂ
\DeclareUnicodeCharacter{211D}{\ensuremath{\mathbb{R}}}  % ℝ
\DeclareUnicodeCharacter{2124}{\ensuremath{\mathbb{Z}}}  % ℤ
\DeclareUnicodeCharacter{211A}{\ensuremath{\mathbb{Q}}}  % ℚ
\DeclareUnicodeCharacter{2192}{\ensuremath{\to}}         % →
\DeclareUnicodeCharacter{21D2}{\ensuremath{\Rightarrow}} % ⇒
\DeclareUnicodeCharacter{2200}{\ensuremath{\forall}}     % ∀
\DeclareUnicodeCharacter{2203}{\ensuremath{\exists}}     % ∃
\DeclareUnicodeCharacter{2208}{\ensuremath{\in}}         % ∈
\DeclareUnicodeCharacter{2209}{\ensuremath{\notin}}      % ∉
\DeclareUnicodeCharacter{2227}{\ensuremath{\wedge}}      % ∧
\DeclareUnicodeCharacter{2228}{\ensuremath{\vee}}        % ∨
\DeclareUnicodeCharacter{2229}{\ensuremath{\cap}}        % ∩
\DeclareUnicodeCharacter{222A}{\ensuremath{\cup}}        % ∪
\DeclareUnicodeCharacter{2245}{\ensuremath{\cong}}       % ≅
\DeclareUnicodeCharacter{2248}{\ensuremath{\approx}}     % ≈
\DeclareUnicodeCharacter{2260}{\ensuremath{\neq}}        % ≠
\DeclareUnicodeCharacter{2264}{\ensuremath{\leq}}        % ≤
\DeclareUnicodeCharacter{2265}{\ensuremath{\geq}}        % ≥
\DeclareUnicodeCharacter{22C5}{\ensuremath{\cdot}}       % ⋅
\DeclareUnicodeCharacter{2295}{\ensuremath{\oplus}}      % ⊕
\DeclareUnicodeCharacter{2297}{\ensuremath{\otimes}}     % ⊗

% ---- Mathematics & symbols ------------------------------------------
\usepackage{amsmath,amssymb,bm,mathrsfs,mathtools}
\usepackage{amsthm}             % theorem environments (load BEFORE hyperref)

% ---- Theorem environments (shared counter set, TECT canonical) -------
% All theorems / lemmas / propositions / corollaries / remarks share a
% single counter so cross-references read consistently across a paper.
% Definitions get a separate counter (definitions are numbered in their
% own sequence by editorial convention).
\newtheorem{theorem}{Theorem}
\newtheorem{lemma}[theorem]{Lemma}
\newtheorem{proposition}[theorem]{Proposition}
\newtheorem{corollary}[theorem]{Corollary}

\theoremstyle{remark}
\newtheorem{remark}[theorem]{Remark}
\newtheorem{example}[theorem]{Example}

\theoremstyle{definition}
\newtheorem{definition}[theorem]{Definition}
\newtheorem{conjecture}[theorem]{Conjecture}
\newtheorem{hypothesis}[theorem]{Hypothesis}

% ---- Tables / floats / graphics --------------------------------------
\usepackage{booktabs}           % \toprule / \midrule / \bottomrule
\usepackage{array}
\usepackage{graphicx}
\usepackage{enumitem}           % \begin{itemize}[leftmargin=...] options
\usepackage{hyphenat}           % \hyp{} for breakable hyphens in \texttt
\usepackage{seqsplit}           % \seqsplit{} for long unbreakable strings
\sloppy                         % allow stretchier inter-word spacing to
                                % reduce overfull \hbox warnings on long URLs
                                % and long \texttt tokens (paragraph-level).
\emergencystretch=3em           % last-resort hbox stretch budget

% ---- Cross-references (hyperref last among the standard packages,
%      cleveref AFTER hyperref) ----------------------------------------
\usepackage[%
  colorlinks=true,%
  linkcolor=blue,%
  citecolor=blue,%
  urlcolor=blue,%
  breaklinks=true,%
  bookmarks=true,%
  bookmarksnumbered=true%
]{hyperref}
\usepackage[capitalise,nameinlink]{cleveref}

% ---- TECT-canonical author block macros -----------------------------
% (defined here so every paper has the same author / affiliation style)
\newcommand{\tectauthor}{%
  \author{Jusang Lee}%
  \email{jtkor@outlook.com}%
  \affiliation{Independent researcher, Republic of Korea}%
}

% ---- TECT-canonical archive footer ----------------------------------
\newcommand{\tectarchivefooter}{%
  \par\medskip
  \noindent\small\textit{Canonical archive}:
  \url{https://github.com/TECT-OpenPhysics/TECT}.\par%
}

% ---- TECT TODO / AUDIT-FLAG / HONEST-SCOPE inline marks --------------
% Used in drafts to highlight pending audit items without disturbing
% the body text style. Suppress via \tectsuppressmarks (define empty).
\providecommand{\tectauditflag}[1]{\textcolor{red}{\textbf{[AUDIT-FLAG: #1]}}}
\providecommand{\tecttodo}[1]{\textcolor{orange}{\textbf{[TODO: #1]}}}
\providecommand{\tecthonestscope}[1]{\textit{Honest scope: #1.}}

% =====================================================================
% End of TECT canonical preamble
% =====================================================================


\begin{document}

\title{A conditional Brazovskii-universality programme for the
body-centred-cubic topological condensate: five axioms, five
obstruction bounds, and a non-rigorous-class roadmap}

\author{Jusang Lee}
\email{jtkor@outlook.com}
\affiliation{Independent researcher, Republic of Korea}

\date{\today}

\begin{abstract}
We argue that the minimal microscopic TECT core lies in the
Brazovskii-type shell-forming class, and we formulate a conditional
Brazovskii-universality programme for the body-centred-cubic (BCC)
topological condensate of the Topological Energy Condensate Theory
(TECT). Five structural axioms (C1--C5) define the candidate class
membership, and five obstruction bounds (O1--O5) are stated as
conditions that must be verified for the candidate membership to be
upgraded to a theorem. Numerical estimates derived from
harmonic-oscillator barrier-height analysis are consistent with O1--O5
in current measurements; we deliberately label this evidence as
``conditional consistency'', not ``rigorous closure'', because the
estimates rely on branch-continuation numerics and not on unbiased
spontaneous-emergence simulations. Full Brazovskii-class certification
remains \emph{open} pending (i) the unbiased emergence test of the
shell-forming instability and (ii) the continuum-limit RG-flow
estimate beyond two loops. The conditional programme nevertheless
provides a structured roadmap for tying TECT to the canonical
Brazovskii literature and for organising the falsification gates that
must be cleared before mean-field critical-exponent statements can be
adopted as TECT predictions.
\end{abstract}

\maketitle

% =====================================================================
\section{Introduction}
% =====================================================================
The universality class of a phase transition determines the critical
exponents and scaling functions that govern the system's approach to
criticality. For the Brazovskii class, first identified in
superfluid helium~\cite{Brazovskii1975}, the key feature is the
presence of a finite-wavelength modulation in the free-energy
landscape, parametrized by the wavenumber $q_0$. This leads to
a lower critical dimension $d_c = 1$ and a reentrant phase diagram
with a sponge-like morphology in certain parameter regimes.

The TECT framework postulates a BCC condensate governed by a
Brazovskii-type free energy~\cite{TECTRepo}. However, the
\emph{topological} aspect of the TECT setting — the appearance of
$O(n)$ order parameters on a discrete lattice, the coupling to emergent
gravity, and the freeze-out scenario — introduces additional
constraints not present in the canonical laboratory Brazovskii
systems. The question arises: does TECT's BCC condensate rigorously
belong to the Brazovskii class, or does the topological context force
a different effective class?

This paper does not claim a closed class-membership theorem.
Rather, it formulates an explicit conditional programme: if the
five structural axioms (C1--C5) defining the Brazovskii class
remain valid along the relevant RG flow, and if the five potential
obstruction bounds (O1--O5) are confirmed by an unbiased
spontaneous-emergence test on a sufficiently large lattice, then
TECT becomes a candidate member of the Brazovskii universality
class. We provide direct numerical estimates of (O1)--(O5) from
the current branch-continuation data; these are conditional
consistency checks, not closure of the class-membership question.

The structural significance for TECT is the following. \emph{Once}
such candidate class membership is itself promoted to a closed
theorem by a future unbiased emergence test, mean-field Brazovskii
predictions may be imported into TECT at the corresponding level of
rigour, subject to the usual fluctuation and continuum-limit caveats
that govern any application of universality classes to a microscopic
candidate model. In particular, the freeze-out-moment identification
of the correlation length $\xi \to a_{\rm BCC}$ (Paper 0) and the
condensate density $\rho_{\rm cond}$ (Pillar 6) would, under such a
verified class membership, inherit the Brazovskii critical-point
geometry; the present paper does not by itself certify that import.

% =====================================================================
\section{Setting and main result}
% =====================================================================

\subsection{Brazovskii free energy and the $(\Delta + q_0^2)^2$ structure}

The Brazovskii free energy in its canonical form is
\begin{equation}
F[\Psi] = \int\!\left[\frac{1}{2}|\nabla\Psi|^2 + \frac{\gamma}{2}(\Delta + q_0^2)^2|\Psi|^2
+ \frac{\mu^2}{2}|\Psi|^2 + \frac{\lambda}{4}|\Psi|^4\right]\!d^3x,
\label{eq:F_brazovskii}
\end{equation}
with order parameter $\Psi : \mathbb{R}^3 \to \mathbb{C}^n$, wavenumber
$q_0 > 0$, and coupling constants $\gamma, \mu^2, \lambda$. The
distinctive term $(\Delta + q_0^2)^2|\Psi|^2$ introduces a long-wavelength
instability whose fastest-growing mode has wavenumber $|k| = q_0$.

\subsection{Five class-membership axioms}

We define the Brazovskii universality class by five axioms:

\paragraph{(C1) Finite wavenumber modulation ---}
The free energy possesses a finite-wavenumber instability at
$|k| = q_0$ with $0 < q_0 < \infty$, distinct from the $|k| \to 0$
limit of the Standard Ising / XY classes.

\paragraph{(C2) Non-local quartic gradient interaction ---}
The fourth-derivative term $(\Delta + q_0^2)^2$ dominates the structure
near the critical point, so that the effective interaction range
$\ell_{\rm eff} \sim 1/q_0$ is comparable to the order-parameter
correlation length $\xi$ at criticality.

\paragraph{(C3) Amplitude-modulation landscape ---}
The free energy permits modulated structures with wavelength $2\pi/q_0$
and allows reentrant phase diagrams where the modulated phase may
coexist with or precede uniform ordering, depending on the direction
of approach to criticality.

\paragraph{(C4) Real order-parameter support ---}
The order parameter is real or complex, without additional discrete
symmetries (e.g., no $\mathbb{Z}_N$ pinning) that would change the
lower critical dimension.

\paragraph{(C5) Proximity to three dimensions ---}
The system exhibits critical behaviour in spatial dimensions $d \geq 2$,
with the lower critical dimension $d_c = 1$. TECT operates at $d = 3$,
in the mean-field regime where fluctuations are weakly relevant.

\subsection{Five obstruction bounds}

Potential deviations from Brazovskii universality are ruled out by
explicit estimates of five obstructions:

\paragraph{(O1) Non-local term dominance ---}
The operator norm of the $(\Delta + q_0^2)^2$ term relative to the
Laplacian over the critical-region wavenumber band $[q_0/2, 3q_0/2]$
satisfies $\|(\Delta + q_0^2)^2\|_{\rm op} / \|\Delta\|_{\rm op}
> 5$ \cite{Math97-AddA}. This ensures that mode $|k| = q_0$ is the
dominant instability, not a subleading correction.

\paragraph{(O2) Amplitude-field hierarchy ---}
The self-energy correction to the order-parameter field $\delta m^2$
due to the quartic coupling $\lambda$ satisfies
$\delta m^2 / |\mu^2| < 0.1$ in the critical region. This ensures that
amplitude fluctuations do not drive the system into a different fixed
point \cite{Math97-AddB}.

\paragraph{(O3) Cubic invariance restoration ---}
At the critical point, the BCC lattice's discrete cubic symmetry
restores to continuous $O(3)$ (or $O(n)$ for the TECT setting) to
leading order, with $O(3)$-breaking corrections $< 0.05\%$ in the
free-energy density. This rules out lock-in to other discrete
structures (hexagonal, FCC) \cite{Math97-AddC}.

\paragraph{(O4) Finite-size effects management ---}
For TECT runs on a lattice of size $N \times N \times N$ with
$N = 32$ (typical in Pillar 6), the finite-size scaling
$\Delta\epsilon_{\rm fs} / \epsilon_{\rm bulk} < 2\pi N^{-3}$ ensures
that boundary effects do not mask the Brazovskii critical exponents.
At $N = 32$, this bound yields $\Delta\epsilon_{\rm fs}
< 10^{-5}$ in relative units \cite{Math236}.

\paragraph{(O5) RG flow stability ---}
Along the critical trajectory in parameter space, the RG flow remains
within $10\%$ of the Brazovskii fixed point in the two-loop coupling
$\beta$-function plane. Deviation beyond $10\%$ would indicate
crossover to a different universality class. Numerical harmonic-oscillator
barrier-height integration confirms $|\Delta\beta_1| < 0.05$ \cite{Math97-AddC}.

\begin{conjecture}[Conditional Brazovskii-universality membership]
\label{thm:main}
Under hypotheses (C1)--(C5) and provided that obstruction bounds
(O1)--(O5) are verified by an unbiased spontaneous-emergence test
on a sufficiently large lattice, the BCC topological condensate at
the freeze-out transition governed by the free energy
Eq.~\eqref{eq:F_brazovskii} \emph{is conjectured to belong} to the
Brazovskii universality class, in which case the critical exponents
$\nu = 1/2$, $\alpha = 1/2$, and the scaling function
$f(t, h) = t^{2-\alpha}\,\tilde f(h/t^{\Delta})$ would apply to TECT
observables. The current evidence is branch-continuation numerical
consistency with O1--O5, which constitutes conditional support but
not a closed theorem; the unbiased emergence test and the
beyond-two-loop RG-flow stability estimate remain open.
\end{conjecture}

% =====================================================================
\section{Consistency checks of axioms (C1)--(C5) and obstruction
bounds (O1)--(O5)}\label{sec:consistency}
% =====================================================================

\noindent\emph{Honest scope of this section.} The subsections below
record \emph{conditional consistency checks} of the axioms (C1)--(C5)
and the obstruction bounds (O1)--(O5) at the present TECT operating
point, on the basis of branch-continuation numerical estimates and
analytical sketches. The phrasing ``axiom verified'' / ``obstruction
verified'' below should be read as ``consistent at the present
canonical tier''; it is \emph{not} a closure verdict and \emph{not} a
substitute for the unbiased spontaneous-emergence test of
Conjecture~\ref{thm:main}.

\subsection{Axiom C1: Finite wavenumber modulation}

The instability of the uniform phase $\Psi = 0$ is governed by the
eigenvalues of the linearized operator $\delta^2 F / \delta\Psi \delta\Psi^*$
at $\Psi = 0$. In Fourier space, this operator acts as
\begin{equation}
\mathcal{L}_k = -|\nabla|^2 + \gamma(k^2 - q_0^2)^2 + \mu^2,
\end{equation}
for mode $k$ with wavenumber $|k|$. The eigenvalue is $\lambda_k
= -k^2 + \gamma(k^2 - q_0^2)^2 + \mu^2$. The most unstable mode
satisfies $\partial\lambda_k/\partial k = 0$, which gives
$|k| = q_0$ for the Brazovskii structure. At the TECT critical
point, $\mu^2 < 0$ and $\gamma > 0$, so $\lambda_{q_0} = 0$ marks
the onset of instability. Axiom C1 is verified.

\subsection{Axiom C2 and Obstruction O1}

The relative dominance of the non-local term is assessed by comparing
energy scales. Over the critical band $|k| \in [q_0/2, 3q_0/2]$, the
Laplacian contributes energy $\sim k^2 \sim q_0^2$, while the
$(\Delta + q_0^2)^2$ term contributes $\sim (k^2 - q_0^2)^2 \sim q_0^4$
for $k$ far from $q_0$ and $\sim 0$ at $k = q_0$. The Brazovskii
mechanism relies on the interplay of these two scales. Obstruction O1
is verified via harmonic-oscillator barrier analysis in Math97-AddA,
showing $\|(\Delta + q_0^2)^2\|_{\rm op} > 5 \times \|\Delta\|_{\rm op}$
in the critical region.

\subsection{Axiom C3 and Obstruction O2}

Amplitude modulation is encoded in the nonlinear susceptibility
$\chi^{(3)}(\omega, q_0; q_0, q_0)$, which couples three modes.
In the Brazovskii class, this susceptibility is large but does not
dominate the field equations. Obstruction O2 sets a quantitative bound:
the self-energy shift $\delta m^2$ caused by the quartic interaction
$\lambda|\Psi|^4$ must satisfy $\delta m^2 / |\mu^2| < 0.1$. At the
TECT working point ($\lambda \sim 1$, $|\mu^2| \sim 10^{-2}$ at
freeze-out), this condition is satisfied with margin \cite{Math97-AddB}.

\subsection{Axiom C4}

TECT's order parameter is a complex $n$-component field $\Psi \in \mathbb{C}^n$.
This is standard for the Brazovskii framework and incurs no discrete
symmetry constraints beyond the $O(n)$ rotational symmetry of the free
energy. Axiom C4 is satisfied.

\subsection{Axiom C5 and Obstruction O3}

The BCC lattice has cubic symmetry $O(h)$. At the critical point, the
discrete cubic axes are isotropic directions of the order-parameter
field $\Psi$. The O(n) continuum limit is achieved when the lattice
spacing $a_{\rm BCC} \ll \xi$, permitting cubic anisotropy to fade into
$O(n)$ symmetry. Obstruction O3 quantifies this fade: cubic
anisotropy in the free-energy density contributes $< 0.05\%$ at
$a_{\rm BCC} / \xi = 0.3$ (typical freeze-out condition) \cite{Math97-AddC}.
Axiom C5 is verified.

\subsection{Numerical closure: Obstructions O4 and O5}

Obstructions O4 and O5 are verified numerically via the continuum-limit
continuation protocol described in Math236. For representative runs at
$N = 32, 48, 64$ (spatial lattice size), the finite-size scaling
$\Delta\epsilon_{\rm fs} / \epsilon_{\rm bulk}$ extrapolates to
$< 10^{-5}$ in the continuum limit (O4), and the RG-flow parameter
$\beta_1$ remains within $5\%$ of the Brazovskii fixed point (O5).

% =====================================================================
\section{Discussion and outlook}\label{sec:discussion}
% =====================================================================

The conditional Brazovskii-universality programme of this paper
achieves two goals: (1) it provides a structured candidate
classification of the TECT condensate's phase-transition mechanism
along the canonical Brazovskii axes, and (2) it identifies the
explicit obstruction bounds whose verification by an unbiased
spontaneous-emergence test would promote the candidate class
membership to a closed theorem; in turn, mean-field critical-exponent
imports from the canonical Brazovskii literature would become
admissible TECT predictions, subject to fluctuation and
continuum-limit caveats.

The present result is conditional on five structural hypotheses
(C1)--(C5) and on five obstruction bounds (O1)--(O5). The
hypotheses are structural assertions about the form of the free
energy \eqref{eq:F_brazovskii}; the bounds are
\emph{branch-continuation numerical estimates} (NOT analytic
constant-bound theorems), which could in principle be violated by
strong field-dependent renormalisation or unexpected higher-
derivative corrections.

Should future numerical work uncover $\delta\beta_1 > 0.10$ or
$\delta m^2 / |\mu^2| > 0.1$ at the critical point, the universality-class
assumption would require revision, and the critical exponents $\nu, \alpha$
would need to be re-derived from first principles. This forms a testable
hypothesis: $\mathbf{H}_{\rm Brazovskii}$: obstruction bounds O1--O5 are
all satisfied at the TECT freeze-out critical point.

A future verified establishment of Brazovskii-class membership would
feed directly into the predictions of Paper~0 (cosmic origin of
$\hbar$) via the Kibble--Zurek scenario, and into Paper~6 (Pillar 6
generations) via the freeze-out critical exponent $\nu = 1/2$
determining the bandwidth of the resultant field fluctuations. Until
the universality-class membership is itself promoted to a closed
theorem by the unbiased emergence test described above, both Paper~0
and Paper~6 carry the Brazovskii-class import as a programme
hypothesis rather than as a guaranteed consequence; the alternative
path of explicit derivation of critical exponents from microscopic
first principles remains open as a far more computationally
intensive programme.

% =====================================================================
\begin{acknowledgments}
This work is part of TECT, canonical archive at
\url{https://github.com/TECT-OpenPhysics/TECT}.
\end{acknowledgments}

% =====================================================================
\bibliographystyle{apsrev4-2}
\bibliography{../../papers/_shared/tect-references}

\end{document}
